\documentclass[a4paper,12pt]{article}
\usepackage[utf8]{inputenc}
\usepackage[T1]{fontenc}
\usepackage[ngerman]{babel}
\usepackage{amsmath}
\usepackage{amssymb}
\usepackage{enumitem}
\usepackage{geometry}
\geometry{a4paper, margin=2.5cm}

\title{Einsendeaufgaben -- Aufgabe 1.4\\[0.5em]
\large 61211 Analysis}
\author{}
\date{}

\begin{document}
\maketitle

\section*{Aufgabe 1.4}

\begin{enumerate}[label=(\roman*)]

    \item
    \underline{1. $\forall f \in X : \|f\|_X \geq 0$}

    Für alle $x \in [0, 1]$ gilt
    \[
        (1 + x e^x)|f(x)| \geq 0.
    \]
    Demnach gilt
    \[
        \|f\|_X = \int_0^1 (1 + x e^x)|f(x)|\, dx \geq 0.
    \]

    \underline{2. $\forall f \in X : \|f\|_X = 0 \Longleftrightarrow f = \hat{0}$}

    \glqq$\Rightarrow$\grqq-Richtung:

    Sei $\|f\|_X = 0$. Da $g(x) = (1 + x e^x)|f(x)|$ stetig ist und das Riemann-Integral
    \[
        \int_0^1 g(x)\, dx = 0
    \]
    ist, muss $f = \hat{0}$ sein.
    Für nicht stetige Funktionen gilt dies nicht allgemein.

    \glqq$\Leftarrow$\grqq-Richtung:

    Sei $f = \hat{0}$. Dann gilt direkt
    \begin{align*}
        0 &= \int_0^1 (1 + x e^x) \cdot 0\, dx \\
          &= \int_0^1 (1 + x e^x)|f(x)|\, dx \\
          &= \|f\|_X
    \end{align*}

    \item
    $\forall f \in X\; \forall a \in \mathbb{R} : \|af\|_X = |a|\,\|f\|_X$
    \begin{align*}
        \|af\|_X &= \int_0^1 (1 + x e^x)|af(x)|\, dx \\
                 &= |a| \int_0^1 (1 + x e^x)|f(x)|\, dx \\
                 &= |a|\,\|f\|_X
    \end{align*}

    \item
    $\forall f, g \in X : \|f + g\|_X \leq \|f\|_X + \|g\|_X$
    \begin{align*}
        \|f + g\|_X &= \int_0^1 (1 + x e^x)|f(x) + g(x)|\, dx \\
                    &\leq \int_0^1 (1 + x e^x)(|f(x)| + |g(x)|)\, dx
                    \quad (\text{da } (1 + x e^x) \geq 0) \\
                    &= \int_0^1 (1 + x e^x)|f(x)|\, dx + \int_0^1 (1 + x e^x)|g(x)|\, dx \\
                    &= \|f\|_X + \|g\|_X
    \end{align*}

\end{enumerate}

\end{document}
